\clearpage\pagestyle{otMain}\setlength{\otBlocEspacement}{2.5mm}\fontsize{11}{15}\selectfont\corrigesChapitre{12}{Droites remarquables et raisonnement}
\begin{blocLivre}{cor-12-1}\begin{corrige}{1}{Reconnaître des droites}\demarche{Lire chaque codage : angle droit, milieu, sommet.}
$(d_1)$ est perpendiculaire à $[BC]$ en son milieu I : c’est la $\res{\text{médiatrice de }[BC]}$. $(d_2)$ passe par C et est perpendiculaire à $(AB)$ : c’est la $\res{\text{hauteur issue de C}}$. $(d_3)$ passe par B et par le milieu J de $[AC]$ : c’est la $\res{\text{médiane issue de B}}$.
\label{fig-cor-12-1}\figureGeo[corrige]{rem-01}
\end{corrige}\end{blocLivre}
\begin{blocLivre}{cor-12-2}\begin{corrige}{2}{Hypothèse et conclusion}\demarche{Échanger hypothèse et conclusion pour la réciproque.}
Hypothèse : ABCD est un rectangle. Conclusion : $AC=BD$. Réciproque : « Si un quadrilatère a des diagonales de même longueur, alors c’est un rectangle. » Elle est $\res{\text{fausse}}$ : le trapèze isocèle EFGH de la figure corrigée a des diagonales de même longueur sans être un rectangle.
\label{fig-cor-12-2}\figureGeo[corrige]{rem-02}
\end{corrige}\end{blocLivre}
\begin{blocLivre}{cor-12-3}\begin{corrige}{3}{Une hauteur et une aire}\demarche{Associer une base à la hauteur perpendiculaire.}
$\mathcal A=\frac{AB\times CH}2=\frac{8\times3}2=\res{12\ \mathrm{cm}^2}$. La hauteur CH est perpendiculaire à la droite $(AB)$ qui porte la base ; le côté AC ne l’est pas nécessairement.
\label{fig-cor-12-3}\figureGeo[corrige]{rem-03}
\end{corrige}\end{blocLivre}
\begin{blocLivre}{cor-12-4}\begin{corrige}{4}{Un point de la médiatrice}\demarche{Utiliser la propriété caractéristique dans les deux sens.}
M est sur la médiatrice, donc il est équidistant de A et B : $\res{MB=4{,}5\ \mathrm{cm}}$. N est équidistant de A et B, donc il appartient à la $\res{\text{médiatrice de }[AB]}$ : on utilise ici la réciproque, vraie elle aussi.
\label{fig-cor-12-4}\figureGeo[corrige]{rem-04}
\end{corrige}\end{blocLivre}
\begin{blocLivre}{cor-12-5}\begin{corrige}{5}{Les angles d’un triangle isocèle}\demarche{Répartir la somme restante en deux angles égaux.}
Les angles à la base, en B et en C, sont égaux. Leur somme vaut $180-74=106^\circ$ ; chacun vaut $106:2=\res{53^\circ}$. On utilise la somme des angles d’un triangle et la symétrie du triangle isocèle.
\label{fig-cor-12-5}\figureGeo[corrige]{rem-05}
\end{corrige}\end{blocLivre}
\begin{blocLivre}{cor-12-6}\begin{corrige}{6}{Un contre-exemple}\demarche{Construire deux triangles de même aire et de côtés différents.}
Il suffit d’un contre-exemple. Le triangle rectangle de côtés de l’angle droit 4 cm et 3 cm a pour aire $\frac{4\times3}2=6$ cm$^2$. Le triangle de base 6 cm et de hauteur 2 cm a aussi pour aire $\frac{6\times2}2=6$ cm$^2$, mais une base de 6 cm, plus longue que tous les côtés du premier, dont le plus long mesure 5 cm. L’affirmation est $\res{\text{fausse}}$.
\label{fig-cor-12-6}\figureGeo[corrige]{rem-06}
\end{corrige}\end{blocLivre}
\begin{blocLivre}{cor-12-7}\begin{corrige}{7}{Une médiane partage l’aire}\demarche{Comparer bases et hauteurs.}
Les triangles ABI et AIC ont des bases de même longueur, $BI=IC$, car I est le milieu de $[BC]$. Ils ont la même hauteur AH, perpendiculaire à $(BC)$. Leurs aires $\frac{BI\times AH}2$ et $\frac{IC\times AH}2$ sont donc égales : chacune vaut $30:2=\res{15\ \mathrm{cm}^2}$.
\label{fig-cor-12-7}\figureGeo[corrige]{rem-07}
\end{corrige}\end{blocLivre}
\begin{blocLivre}{cor-12-8}\begin{corrige}{8}{Utiliser la contraposée}\demarche{Nier la conclusion pour nier l’hypothèse.}
P n’est pas équidistant de A et de B, car $5\ne6$. Par la $\res{\text{contraposée}}$ de l’énoncé : si un point n’est pas équidistant de A et de B, alors il n’appartient pas à la médiatrice de $[AB]$. Donc $\res{P\text{ n’est pas sur la médiatrice}}$.
\label{fig-cor-12-8}\figureGeo[corrige]{rem-08}
\end{corrige}\end{blocLivre}
\begin{blocLivre}{cor-12-9}\begin{corrige}{9}{Une réciproque fausse}\demarche{Chercher un quadrilatère qui vérifie la nouvelle hypothèse sans la conclusion.}
Réciproque : « Si un quadrilatère a des diagonales perpendiculaires, alors c’est un losange. » Elle est $\res{\text{fausse}}$ : le cerf-volant EFGH de la figure corrigée a des diagonales perpendiculaires, mais $EF=EH\ne GF$ ; ce n’est pas un losange.
\label{fig-cor-12-9}\figureGeo[corrige]{rem-09}
\end{corrige}\end{blocLivre}
\begin{blocLivre}{cor-12-10}\begin{corrige}{10}{Le triangle isocèle}\demarche{Utiliser la réciproque de la propriété de la médiatrice.}
$AB=AC$ : A est équidistant de B et C, donc il est sur la médiatrice de $[BC]$. Cette médiatrice passe aussi par I, milieu de $[BC]$ : c’est la droite $(AI)$. Ainsi $(AI)$ est la $\res{\text{médiane}}$ issue de A, et elle est perpendiculaire à $(BC)$ : c’est aussi la $\res{\text{hauteur}}$ issue de A.
\label{fig-cor-12-10}\figureGeo[corrige]{rem-10}
\end{corrige}\end{blocLivre}
\begin{blocLivre}{cor-12-11}\begin{corrige}{11}{Le point commun des hauteurs}\demarche{Utiliser la concourance des hauteurs.}
Les trois hauteurs d’un triangle sont concourantes : la hauteur issue de C passe donc par H, et cette hauteur est $(CH)$. Ainsi $\res{(CH)\perp(AB)}$. Si ABC est rectangle en C, $(CA)$ et $(CB)$ sont des hauteurs qui se coupent en C : $\res{H=C}$.
\label{fig-cor-12-11}\figureGeo[corrige]{rem-11}
\end{corrige}\end{blocLivre}
\begin{blocLivre}{cor-12-12}\begin{corrige}{12}{Raisonner par l’absurde}\demarche{Supposer le contraire et chercher une contradiction.}
Supposons qu’un triangle ABC ait deux angles droits, en B et en C. La somme $\widehat B+\widehat C$ vaudrait déjà $180^\circ$, et l’angle en A serait nul : impossible pour un triangle. On peut aussi voir que $(BA)$ et $(CA)$ seraient deux perpendiculaires à $(BC)$, donc parallèles, et ne pourraient pas se couper en A. $\res{\text{L’hypothèse est fausse.}}$
\label{fig-cor-12-12}\figureGeo[corrige]{rem-12}
\end{corrige}\end{blocLivre}
\begin{blocLivre}{cor-12-13}\begin{corrige}{13}{Trois maisons et un puits}\demarche{Enchaîner deux propriétés de la médiatrice.}
On trace les médiatrices de $[AB]$ et de $[BC]$ ; elles se coupent en O. O est sur la médiatrice de $[AB]$, donc $OA=OB$ ; il est sur celle de $[BC]$, donc $OB=OC$. Ainsi $\res{OA=OB=OC}$ : O est le centre du cercle circonscrit à ABC et il appartient aussi à la médiatrice de $[AC]$.
\label{fig-cor-12-13}\figureGeo[corrige]{rem-13}
\end{corrige}\end{blocLivre}
\begin{blocLivre}{cor-12-14}\begin{corrige}{14}{Vrai ou faux ?}\demarche{Démontrer, réfuter ou reconnaître une contraposée.}
a) $\res{\text{Vrai}}$ : les angles correspondants formés avec la droite commune sont droits, donc égaux ; les droites sont parallèles. b) $\res{\text{Faux}}$ : un triangle isocèle d’angles $74^\circ$, $53^\circ$ et $53^\circ$ est un contre-exemple. c) $\res{\text{Vrai}}$ : c’est la contraposée de l’énoncé a), qui est vrai.
\label{fig-cor-12-14}\figureGeo[corrige]{rem-14}
\end{corrige}\end{blocLivre}
\begin{blocLivre}{cor-12-15}\begin{corrige}{15}{Une médiane dans une médiane}\demarche{Utiliser deux fois la propriété de la médiane.}
$[AI]$ est une médiane de ABC : l’aire de ABI vaut $\frac12\times36=18$ cm$^2$. Dans le triangle ABI, $[BJ]$ est une médiane car J est le milieu de $[AI]$ : l’aire de ABJ vaut $\frac12\times18=\res{9\ \mathrm{cm}^2}$, soit $\frac14$ de l’aire de ABC.
\label{fig-cor-12-15}\figureGeo[corrige]{rem-15}
\end{corrige}\end{blocLivre}
